\section{Delayed energy and two-sided kernel information}
\label{lifts:three-groups}

The previous arguments controlled endpoint laws or coset masses. A
kernel on complete decks can instead retain the sets of positions
occupied by three groups of labels and Fourier transform only the
internal permutations of each group. Its remaining matrix has rows
and columns indexed by those position sets. We therefore need both
forward and stationary-reverse marginal estimates. The scalar estimate
below is deliberately valid one layer before the block endpoint:
that extra time statement will justify cancellation in the all-sign
matrix block. The full-deck theorem at the end illustrates how this
two-sided kernel information can be amplified; the additional content
here is the control of its quotient and sign matrices.

\begin{lemma}[A delayed energy contraction]
\label{lifts:delayed-marginal}
For any periodic repetition of a permutation of the $d$ coordinate directions, every
specified ordered tuple of $k\le3n/4$ distinct cards, from every
deterministic start, has total-variation distance at most $n^{-10}$
from the uniform injective tuple after $t\ge2000d-1$ layers, for
sufficiently large $d$. The same bound holds for the stationary
reverse of the physical Thorp shuffle.
\end{lemma}

\begin{proof}
Fix $q$ observed cards and one tag, with $h=n-q\ge n/4$ available
positions. Let $B_t$ be this available set at time $t$. The conditional
tag law $p_t$ on $B_t$ is the path-prefix
flow of Lemma~\ref{lem:marginal-averaging}. Its centered energy
\[
 E_t=\sum_{x\in B_t}(p_t(x)-1/h)^2
\]
is at most one and nonincreasing. In the notation of
Section~\ref{sec:first-route}, $B_t$ is $V_t$, $h$ is $m$, and
$E_t=\|f_t\|_2^2$ is the realized energy. Here layer $s$ means the step from
$s-1$ to $s$; write $j_s=i_{s-1}$ for its coordinate direction.
Let $\mathcal F_s$ contain all switches through layer $s$. Conditioning on this larger past fixes
the path-defined function $p_s$; it does not redefine it as the
tag's law conditional on all card positions.

For $t>d$, put $u=t-d$, $i=j_t=j_u$, and
$H_b=\{x\in V:x_i=b\}$ for $b=0,1$. Between layers $u$
and $t$, every coordinate other than $i$ occurs once. The two
coordinate-$i$ half counts $h_b=|B_u\cap H_b|$ are consequently fixed by
$\mathcal F_u$. The delayed half-space symmetry
\cite[Lemma~\ref*{f-shear:lem:delayed-symmetry}]{coordinateFrames}
states that the terminal pair consisting of the available set and
its centered mass vector is invariant, conditional on
$\mathcal F_u$, under every map $T_zx=x+x_i z$, $z_i=0$.
It applies to every nonempty starting available set and every probability
vector supported on it. The future switches in the intervening distinct coordinates are independent of
$\mathcal F_u$.

For clarity, its energy consequence has no hidden normalization.
Averaging over a uniform $z$ aligns each pair of opposite-half
available positions with probability $2/n$. Set
$a(x)=p_{t-1}(x)-\mathbf1_{B_{t-1}}(x)/h$ for $x\in V$;
the mean loss in the final layer is thus
\[
 \frac1n\sum_{x\in B_{t-1}\cap H_0,\ y\in B_{t-1}\cap H_1}
       (a(x)-a(y))^2
 \ge\frac{\min(h_0,h_1)}n E_{t-1}.
\]
The inequality follows by writing $A=\sum_{B_{t-1}\cap H_0}a$:
the double sum is
$h_1\sum_{H_0}a^2+h_0\sum_{H_1}a^2+2A^2$, since
$\sum_xa(x)=0$. Therefore
\begin{equation}\label{lifts:conditional-shear-loss}
 \E(E_{t-1}-E_t\mid\mathcal F_u)
 \ge\frac{\min(h_0,h_1)}n\E(E_{t-1}\mid\mathcal F_u).
\end{equation}

Call the split bad if $\min(h_0,h_1)<h/4$. Conditional on
$\mathcal F_{u-1}$, its count difference after layer $u$ is a
sum of independent fair signs, one for each pair with exactly one
available position. Its second moment is at most $h$, so the bad
probability is at most $4/h\le16/n$. Crucially, we charge this
event against the earlier energy $E_{u-1}$, which is known before
those signs are drawn, rather than treating it as independent of
$E_{t-1}$. Monotonicity gives, with $a_t=\E E_t$ and $\eta=1/16$,
\begin{equation}\label{lifts:delay-recurrence}
 a_t\le(1-\eta)a_{t-1}+\eta(16/n)a_{t-d-1}\qquad(t>d).
\end{equation}
Set $\gamma=1/64$. The bound $a_t\le e^{-\gamma(t-d)}$ is
trivial for $t\le d$, and induction in the recurrence proves it
thereafter because
\[
 (1-\eta)e^\gamma+(16\eta/n)e^{\gamma(d+1)}\le1
\]
for all sufficiently large $d$. The first term is strictly below
one, and the second tends to zero since $\gamma<\log2$.

The expected tag error is at most $\frac12\sqrt{ha_t}$.
Lemma~\ref{lem:sequential-tv}, after forgetting paths but retaining
endpoints, bounds the tuple error by
\[
 \frac12 k\sqrt{ne^{-\gamma(t-d)}}
 \le\frac12n^{3/2}e^{-(1999d-1)/128}\le n^{-10}.
\]
The last inequality uses $t\ge2000d-1$. A reverse physical step
is the inverse random position map. Undoing its inverse coordinate
rotations again gives a cyclic schedule of coordinate matchings,
with a deterministic change of phase. The proof covers every such
order and phase, proving the reverse assertion.
\end{proof}

\begin{theorem}[Three-group amplification]
\label{lifts:three-group-theorem}
The worst-start full-deck total-variation distance tends to zero after
$30000d$ physical Thorp shuffles. The proof uses fifteen independent
blocks of length $2000d$ and three groups of labels.
\end{theorem}

\begin{proof}
Partition the labels into three groups of sizes $m_i$ differing by
at most one. For large $d$, each $m_i\ge n/4$ and $m_i\to\infty$.
Encode a deck as $(s,g)$, where $s=(S_1,S_2,S_3)$ records the
position sets occupied by the groups. Fix an ordering of each label
group and, for each $s$, a reference bijection
$r_{s,i}:[m_i]\to S_i$. In the deck $(s,g)$, the $a$th label of
group $i$ occupies $r_{s,i}(g_i(a))$, with
$g=(g_1,g_2,g_3)$. Let
\[
 H=\prod_{i=1}^3 S_{m_i},\quad
 \mathcal S=\{\text{ordered position partitions of these sizes}\},
 \quad M=|\mathcal S|=n!/\prod_i m_i!.
\]
The $T=2000d$ kernel $K$ has convolution fibers
\[
 K((s,g),(s',g'))=p_{s,s'}(g'g^{-1}).
\]
Indeed, a position increment $\tau$ taking $s$ to $s'$ changes
$g_i$ to $h_i g_i$, where
$h_i=r_{s',i}^{-1}\tau r_{s,i}$. The law of $(s',h)$ given $s$
is independent of the internal assignment $g$, proving the formula.
The kernel is bistochastic because every position increment acts
bijectively on decks. The reverse fibers are
$\bar p_{s',s}(h)=p_{s,s'}(h^{-1})$.

Put $H_{-i}=\prod_{j\ne i}S_{m_j}$, $M_i=M|H_{-i}|$, and
\[
 p^{(-i)}_{s,s'}(h_{-i})=\sum_{h_i}p_{s,s'}(h).
\]
The positions of labels outside group $i$ determine $(s',h_{-i})$,
whose uniform law has $M_i$ atoms. With $\delta=n^{-10}$,
Lemma~\ref{lifts:delayed-marginal} at $T$ gives absolute row error
at most $2\delta$ from this uniform quotient. Its reverse assertion
gives absolute column error at most $2\delta$, after inversion of
the internal group elements. These are the two quotient hypotheses
of~\cite[Theorem~\ref*{l-l:three-group}]{partialFourier}.

Its remaining scalar hypothesis concerns the product of the three
sign characters. We verify it for the untrimmed kernel. Condition
on the first $T-1$ layers. If two cards of group 1 occupy a pair
in the final matching, flipping that pair's coin preserves the
output position partition $s'$ and changes the product of the
internal signs. The conditional signed contribution to each $s'$
therefore cancels. The absolute row sum of the all-sign block is
at most the probability that no such pair exists. At $T-1$,
Lemma~\ref{lifts:delayed-marginal} makes the group-1 position set
within $\delta$ of a uniform $m_1$-subset. A uniform subset avoids
containing a whole pair with probability
\[
 \prod_{j=0}^{m_1-1}\left(1-\frac{j}{n-j}\right)
 \le\exp\left(-\frac{m_1(m_1-1)}{2n}\right).
\]
This follows by sampling its points in order: after $j$ points with
no completed pair, exactly $j$ partners are forbidden among the
$n-j$ remaining positions. The bound tends to zero. Apply the
same argument to the reverse kernel, using its $T-1$ preceding
reverse layers and a canceling final reverse layer. In physical
coordinates the final reverse step first rotates the coordinates and
then switches a matching. That fixed rotation preserves the uniform
subset law and its variation bound, so the same pair-avoidance estimate
applies just before the canceling switches. This gives the
same bound on the original absolute column sums. Thus the untrimmed
all-sign row and column norms are $o(1)$, as required. A theorem
only at time $T$ would not justify this step.

All hypotheses of the three-group transfer are now verified. More
explicitly, it retains only fibers satisfying
\[
 p^{(-i)}_{s,s'}(h_{-i})\le2/M_i\qquad(i=1,2,3),
\]
losing at most $6\delta$ in each row and column. For each product
type of degrees $D_i$, its retained block $B_\rho$ satisfies
$\|B_\rho\|_{\HS}^2\le4D_i/\prod_{j\ne i}D_j$.
Writing $D=\prod_iD_i$, fifteen factors, including the regular
multiplicity $D$, contribute at most $2^{30}D^{-4}$ for $D>1$.
The transfer also treats all eight scalar types: quotient row and
column bounds handle those with a trivial factor, and the verified
all-sign cancellation handles the last one. It concludes average
row convergence of $K^{15}$ to uniform.

Finally the original shuffle kernel is equivariant under arbitrary
relabelings, and those relabelings act transitively on decks. Every
row of $K^{15}$ consequently has the same distance from uniform,
so average convergence is worst-start convergence. Fifteen blocks
take $15T=30000d$ physical shuffles.
\end{proof}
